% GENERATED by scripts/build-lecture-slides.py; do not edit. Portable profile: XeTeX, theme and fonts from the core.
\makeatletter\IfFileExists{core/beamer/beamerthememx.sty}{\def\mxcorepath{core/}\def\input@path{{core/beamer/}{core/}}}{\def\mxcorepath{../core/}\def\input@path{{../core/beamer/}{../core/}}}\makeatother
\documentclass[aspectratio=169,10pt,t]{beamer}
\usetheme{mx}
\usepackage{graphicx,booktabs}
\usepackage[french]{babel}
\setlength{\parskip}{.45em}
\hypersetup{pdfauthor={},pdfsubject={Lecture slides},pdfkeywords={},pdftitle={Suites numériques}}
\providecommand{\N}{\mathbb{N}}
\providecommand{\Z}{\mathbb{Z}}
\providecommand{\D}{\mathbb{D}}
\providecommand{\Q}{\mathbb{Q}}
\providecommand{\R}{\mathbb{R}}
\providecommand{\C}{\mathbb{C}}
\providecommand{\vect}[1]{\overrightarrow{#1}}
\providecommand{\vi}{\vec{\imath}}
\providecommand{\vj}{\vec{\jmath}}
\providecommand{\vk}{\vec{k}}
\providecommand{\norm}[1]{\left\lVert #1\right\rVert}
\providecommand{\vc}[2]{\begin{pmatrix}#1\\#2\end{pmatrix}}
\providecommand{\deter}[2]{\det\left(#1,#2\right)}
\providecommand{\abs}[1]{\left\lvert #1\right\rvert}
\providecommand{\dd}{\,\mathrm{d}}
\providecommand{\e}{\mathrm{e}}
\providecommand{\Card}{\operatorname{Card}}
\providecommand{\E}{\operatorname{E}}
\providecommand{\V}{\operatorname{V}}
\providecommand{\degre}{^{\circ}}
\providecommand{\un}[1]{\,\mathrm{#1}}
\providecommand{\sci}[2]{#1\times10^{#2}}
\providecommand{\pgcd}{\operatorname{PGCD}}
\providecommand{\ppcm}{\operatorname{PPCM}}
\title{Suites numériques}
\subtitle{Cours et exemples}
\date{}
\author{}
\begin{document}
\mxtitleframe
\begin{frame}{Raisonnement par récurrence}
Soit \(\mathcal{P}_n\) une propriété dépendant d'un paramètre entier \(n\).
            \(\mathcal{P}_n\) est vraie pour tout entier naturel \(n \geq n_0\), \(n_0 \in \mathbb{N}\), si : 

\begin{enumerate}
\setcounter{enumi}{0}
                        \item \(\mathcal{P}_{n_0}\) est vraie (c'est l'\emph{initialisation}).
                        
\item \(\mathcal{P}_k \Rightarrow \mathcal{P}_{k+1}\) pour \(k\geq n_0\) (c'est l'\emph{hérédité}).

\end{enumerate}
\end{frame}
\begin{frame}{Raisonnement par récurrence}
Nous allons démontrer que pour tout entier naturel non nul \(n\),
            la propriété suivante est vraie : 
                  \[
                        \mathcal{P}_n \ : \ 
                        1 + 2 + 3 + \cdots + (n-1) + n = \dfrac{n (n+1)}{2}
                        \;.
                  \]
\end{frame}
\begin{frame}{Raisonnement par récurrence (2)}
\begin{itemize}
\item \textbf{Initialisation :} 
                        Nous avons :
                              \(
                                    \dfrac{1 \cdot (1+1)}{2} = 1 
                              \)
                        ce qui montre \(\mathcal{P}_1\).

\end{itemize}
\end{frame}
\begin{frame}{Raisonnement par récurrence (3)}
\begin{itemize}
\item \textbf{Hérédité :} 
                        \textbf{H}ypothèse de \textbf{R}écurrence :
                        On suppose que, pour un entier \(k \geq 1\) quelconque, \(\mathcal{P}_k\) est vraie.
                        Ceci se traduit alors par 
                              \[
                              1+2+3+\cdots+(k-1)+k = \dfrac{k(k+1)}{2}\;.
                              \]
                        À partir de cette \textbf{HR}, prouvons que \(\mathcal{P}_{k+1}\) est vraie.
                        Pour cela, regardons ce qui se passe au rang suivant (c'est-à-dire au rang \(k+1\)) :
                              \[\begin{aligned}
                              1+2+3+\cdots+(k-1)+k+(k+1) & = \dfrac{k(k+1)}{2}  + (k+1)                      \\
                                                         & = (k+1)\left( \dfrac{k}{2} + 1 \right)            \\
                                                         & = (k+1) \dfrac{k+2}{2}                            \\
                                                         & = \dfrac{(\boldsymbol{k+1}) (\boldsymbol{k+1}+1)}{2}  
                              \end{aligned}\]
                              ce qui montre que \(\mathcal{P}_{k+1}\) est vraie.
                        Ainsi, on vient de prouver que \(\mathcal{P}_k \Rightarrow \mathcal{P}_{k+1}\).

\end{itemize}
\end{frame}
\begin{frame}{Raisonnement par récurrence (4)}
On conclut que la propriété \(\mathcal{P}_n\) est vraie pour tout entier \(n \geq 1\).
\end{frame}
\begin{frame}{Raisonnement par récurrence}
Notons que l'initialisation est nécessaire.
Pour montrer cela, nous allons omettre l'initialisation et montrer
la propriété erronée que \(2^n < 0\) pour tout \(n \in \mathbb{N}\). 

Pour \(n \in \mathbb{N}\), soit \(\mathcal{P}_n\) la propriété << \(2^n\) est négatif >>.
On veut prouver que \(\mathcal{P}_n\) est vraie pour tout \(n \in \mathbb{N}\).
\end{frame}
\begin{frame}{Raisonnement par récurrence (2)}
\begin{enumerate}
\setcounter{enumi}{0}
\item On veut omettre la première étape.

\item Supposons que pour un certain \(k \in \mathbb{N}\) on ait \(2^k < 0\).
      On veut donc montrer \(\mathcal{P}_{k+1}\), c'est-à-dire que \(2^{k+1}<0\).
      D'après l'hypothèse de récurrence, on a \(2^k < 0\) et, en multipliant l'inégalité par \(2\), on a :
      \(2^{k+1} < 0\).

\end{enumerate}
\end{frame}
\begin{frame}{Raisonnement par récurrence (3)}
L'hérédité est donc démontrée, et pourtant la propriété est fausse pour tout \(n\) : \(2^n > 0\).
C'est l'initialisation qui manque (\(\mathcal{P}_0\) : << \(2^0 = 1 < 0\) >> est fausse).
\end{frame}
\begin{frame}{Rappel}
Une suite \((U_n)_{n \in E}\) est dite \emph{arithmétique} s'il existe un réel \(r\) tel que \(U_{n+1} = U_n + r\) pour tout \(n \in E\).
Le nombre \(r\) est appelé \emph{raison} de la suite \((U_n)\).
\end{frame}
\begin{frame}{Rappel}
\begin{itemize}
 
\item La suite des nombres impairs, définie par : 
\(
\left\lbrace
\begin{array}{l}
U_0 = 1 \\
U_{n+1} = U_n + 2
\end{array}
\right.
\)
est une suite arithmétique.

\item La suite des nombres entiers naturels est la suite arithmétique de premier terme \(U_0 = 0\) et de raison \(r = 1\).

\end{itemize}
\end{frame}
\begin{frame}{Rappel}
Expression en fonction de \(n\)

 Soit \((U_n)\) une suite arithmétique de premier terme \(U_0\) et de raison \(r\). Alors 
\(U_n = U_0 + n \; r\).
\end{frame}
\begin{frame}{Rappel}
Soit \((U_n)\) une suite arithmétique de raison \(r\).
Si \(p\) et \(q\) sont deux entiers naturels tels que \(p \leq q\), alors : \(U_q = U_p + (q - p) \; r \).
\end{frame}
\begin{frame}{Rappel}
Soit \((U_n)\) la suite arithmétique de raison \(3\) et de premier terme \(U_0 = 1\).
Alors \(U_{20} = U_0 + (20 - 0) \; 3 = 1 + 60 = 61\).
\end{frame}
\begin{frame}{Rappel}
\textbf{Somme des termes d'une suite arithmétique}

 Soit \((U_n)\) une suite arithmétique de raison \(r\).
Posons \(S_n = U_0 + U_1 + \cdots + U_n\).
On a alors
\(\displaystyle{ S_n = \dfrac{(U_0 + U_n) \cdot (n+1)}{2} }\).

				\(
					S = \dfrac{(
					U_0
					+
					U_n
					)
					\cdot
					(
					n+1
					)
					}{
					2
					}
					=
					\dfrac{ (\text{Premier terme} 
					+ 
					\text{Dernier terme} )
					\cdot
					( \text{Nombre de termes} )
					}{ 2 }
				\)
\end{frame}
\begin{frame}{Rappel}
Calculer la somme des entiers naturels de \(1\) à \(100\).
Posons \(U_n = n\) pour \(n \in \mathbb{N}\) : c'est la suite arithmétique de premier terme \(U_0 = 0\) et de raison \(1\).
Comme \(U_0 = 0\), on a \(1 + 2 + \cdots + 100 = U_0 + U_1 + \cdots + U_{100} = S_{100}\), et la formule donne
\[
S_{100} = \dfrac{(U_0 + U_{100}) \cdot (100+1)}{2} = \dfrac{(0 + 100) \cdot 101}{2} = 5050 \; .
\]
On peut aussi utiliser la forme << (premier terme \(+\) dernier terme) \(\times\) nombre de termes \(/ \, 2\) >> sur les \(100\) termes
\(U_1, \ldots, U_{100}\) : \(\dfrac{(1 + 100) \cdot 100}{2} = 5050\).
\end{frame}
\begin{frame}{Rappel (2)}
Calculer la somme des \(n + 1\) premiers nombres impairs : 

Posons \(V_n = 2n+1\) et \(T_n = V_0 + V_1 + \cdots + V_n\).
\[\begin{aligned}
T_n & = \dfrac{(1 + (2n+1)) \cdot (n+1)}{2}       
      = \dfrac{(2 n + 2) \cdot (n+1)}{2}          \\
    & = \dfrac{\not 2 (n+1) (n+1)}{\not 2}        
      = (n+1)^2                                  \;.
\end{aligned}\]
\end{frame}
\begin{frame}{Rappel}
Une suite \((U_n)_{n \in E}\) est dite \emph{géométrique} s'il existe un réel \(q\) tel que \(U_{n+1} = q \cdot U_n\) pour tout \(n \in E\).
Le nombre \(q\) est appelé \emph{raison} de la suite \((U_n)\).
\end{frame}
\begin{frame}{Rappel}
\begin{itemize}
 
\item La suite définie par : 
\(
\left\lbrace
\begin{array}{l}
U_0 = 1/3 \\
U_{n+1} = 2 \cdot U_n
\end{array}
\right.
\)
est une suite géométrique.

\item La suite des puissances de \(2\) est une suite géométrique de raison \(2\).

\end{itemize}
\end{frame}
\begin{frame}{Rappel}
Expression en fonction de \(n\)

 Soit \((U_n)\) une suite géométrique de premier terme \(U_0\) et de raison \(q\). Alors
\(U_n = U_0 \cdot q^n \).
\end{frame}
\begin{frame}{Rappel}
Soit \((U_n)\) une suite géométrique de raison \(q\).
Si \(r\) et \(s\) sont deux entiers naturels tels que \(r \leq s\), alors : \(U_s = U_r \; q^{s-r} \).
\end{frame}
\begin{frame}{Rappel}
Soit \((U_n)\) la suite géométrique de raison \(1/2\) et de premier terme \(U_0 = 8\).
Alors \(U_{12} = U_0\cdot \left(\dfrac{1}{2}\right)^{12 - 0} = 8 \; \dfrac{1}{2^{12}} = \dfrac{1}{512}\).
\end{frame}
\begin{frame}{Rappel}
\textbf{Somme des termes d'une suite géométrique}

 Soit \((U_n)\) une suite géométrique de raison \(q\).
Posons \(S_n = U_0 + U_1 + \cdots + U_n\).
Si \(q \neq 1\), on a
\(\displaystyle{ S_n = U_0 \dfrac{1 - q^{n+1}}{1-q} }\).
	
					\(
						S_n = U_0
						\cdot
						\dfrac{1-q^{n+1}}
						{1-q}
						= \text{\text{Premier terme}} \cdot 
						\dfrac{1-\text{Raison}^{ 						{\text{\footnotesize{Nombre de termes}}}; }}
						{ 1-\text{Raison} }
					\)

Si \(q = 1\), alors \(S_n = (n+1)U_0\).
\end{frame}
\begin{frame}{Rappel}
Calculer la somme \(S = 1 + 2^1 + 2^2 + \cdots + 2^{10}\).
On a 
\(
S = 1 \cdot \dfrac{1 - 2^{11}}{1- 2} = \dfrac{1 - 2^{11}}{-1} = 2^{11} - 1 = 2047.
\)
\end{frame}
\begin{frame}{Suites numériques}
Une suite numérique \((U_n)_{n \in E}\) est une application de \(E \subseteq \mathbb{N}\) vers \(\mathbb{R}\).
\end{frame}
\begin{frame}{Suites numériques}
Terminologies : 

\begin{enumerate}
\setcounter{enumi}{0}
\item La suite est notée \((U_n)\).

\item Le terme de rang \(n\) est noté \(U_n\).

\end{enumerate}
\end{frame}
\begin{frame}{Suites numériques}
Définie par une relation : \(U_{n} = f(n)\)

 On peut définir une suite par la donnée de l'expression du terme de rang \(n\) en fonction de \(n\).
\end{frame}
\begin{frame}{Suites numériques}
Soit \(U_n = 2n + 1\).
      La suite \((U_n)\) est ici défini pour tout \(n \in \mathbb{N}\).
      \[
            U_0 = 1 \qquad U_1 = 3 \qquad U_2 = 5 \qquad U_3 = 7 \; .
      \]
\end{frame}
\begin{frame}{Suites numériques}
Définie par une relation : \(U_{n+1} = f(U_n)\)

 On peut aussi définir une suite par la donnée d'un terme et d'une relation de récurrence.
\end{frame}
\begin{frame}{Suites numériques}
Soit \((U_n)\) la suite définie par : 
      \[
            \left\lbrace
                  \begin{matrix}
                        U_0 = 1 \\
                        U_{n+1} = U_n + 2
                  \end{matrix}
            \right.
      \]
      Les premier termes de cette suite sont : 
      \[
            U_0 = 1 \qquad U_1 = 3 \qquad U_2 = 5 \qquad U_3 = 7 \; .
      \]
\end{frame}
\begin{frame}{Suites numériques}
\begin{itemize}
 
\item Une suite \((U_n)\) est dite \emph{croissante} si pour tout \(n \in \mathbb{N}\) on a : \(U_{n} \leq U_{n+1}\).

\item Une suite \((U_n)\) est dite \emph{décroissante} si pour tout \(n \in \mathbb{N}\) on a : \(U_{n+1} \leq U_{n}\).

\item Une suite est dite \emph{monotone} si elle est croissante ou décroissante.

\end{itemize}
\end{frame}
\begin{frame}{Suites numériques}
\begin{itemize}
 
\item Une suite \((U_n)\) est dite \emph{majorée} s'il éxiste \(M \in \mathbb{R}\) tel que \(U_n \leq M\) pour tout \(n \in \mathbb{N}\).

\item Une suite \((U_n)\) est dite \emph{minorée} s'il éxiste \(m \in \mathbb{R}\) tel que \(U_n \geq m\) pour tout \(n \in \mathbb{N}\).

\item Une suite est dite \emph{bornée} si elle est à la fois minorée et majorée.

\end{itemize}
\end{frame}
\begin{frame}{Suites numériques}
Soit \((U_n)\) une suite numérique et soit \(\ell \in \mathbb{R}\).
On dit que \((U_n)\) converge vers \(\ell\) si tout intervalle ouvert contenant \(\ell\) contient tous les termes \(U_n\)
à partir d'un certain rang. On note alors
\[
\lim_{n \rightarrow + \infty} U_n = \ell \; .
\]
Si elle existe, cette limite est unique.
Une suite qui n'est pas convergente est dite \emph{divergente}.
\end{frame}
\begin{frame}{Suites numériques}
Soient \(f\) une fonction définie sur \(\mathbb{R}^+\) et \((U_n)\) la suite définie par \({U_n = f(n)}\).
Soit \(\ell\) un réel, ou \(+\infty\), ou \(-\infty\).
Alors, 
\[
\text{Si } \lim_{x \rightarrow + \infty} f(x) = \ell \text{ alors } \lim_{n \rightarrow + \infty} U_n = \ell      \; .
\]
\end{frame}
\begin{frame}{Suites numériques}
Soit \((U_n)\) une suite numérique :

\begin{enumerate}
\setcounter{enumi}{0}
\item Si \((U_n)\) est croissante et majorée alors \((U_n)\) est convergente.

\item Si \((U_n)\) est décroissante et minorée alors \((U_n)\) est convergente.

\end{enumerate}
\end{frame}
\begin{frame}{Suites numériques}
Considérons les suites numériques suivantes:
\end{frame}
\begin{frame}{Suites numériques (2)}
\begin{enumerate}
\setcounter{enumi}{0}\item Soit \((W_n)_{n \in \mathbb{N}^*}\) la suite définie par \(W_n = \dfrac{1}{n}\). 
      Comme \(W_n \geq 0\) pour tout \(n\) et que \((W_n)\) est décroissante, alors la suite \((W_n)\) est convergente.

\end{enumerate}
\end{frame}
\begin{frame}{Suites numériques (3)}
\begin{enumerate}
\setcounter{enumi}{1}\item L'exemple précédent est un peu trop facile. Considérons la suite \((X_n)_{n \in \mathbb{N}}\) définie par 
      \(X_n = \sum_{k = 1}^n \dfrac{1}{k \cdot 2^k}\).
      Comme \(X_{n+1} - X_n = \dfrac{1}{(n+1) \cdot 2^{n+1}} > 0\), la suite \((X_n)\) est croissante.
      De plus, comme \(\dfrac{1}{k \cdot 2^k} \leq \dfrac{1}{2^k}\) alors 
      \(X_n \leq \sum_{k = 1}^n \dfrac{1}{2^k}\) pour tout \(n \in \mathbb{N}^*\).
      Or \(\sum_{k = 1}^n \dfrac{1}{2^k} = \dfrac{1}{2} \cdot \dfrac{1 - (1/2)^n}{1 - 1/2} = 1 - \dfrac{1}{2^n} \leq 1\) ; donc \(X_n \leq 1\).
      La suite \((X_n)\) est ainsi majorée; comme elle est croissante alors elle est convergente.

\end{enumerate}
\end{frame}
\begin{frame}{Suites numériques}
Soient \((U_n)\) une suite numérique et \(f\) une fonction numérique.
Supposons que \((U_n)\) converge vers un nombre réel \(\ell\) et que \(f\) est continue dans un voisinage de \(\ell\).
Alors la suite \((f(U_n))\) converge vers \(f(\ell)\).
\end{frame}
\begin{frame}{Suites numériques}
On veut étudier la convergence de la suite \((V_n)\) définie pour tout \(n \in \mathbb{N}^*\) par \(V_n = \sin \dfrac{1}{n}\).
Considérons la suite \((U_n)\) définie par \(U_n = \dfrac{1}{n}\) et la fonction \(f\) définie par \(f(x) = \sin x\).
On sait que la fonction sinus est continue sur \(\mathbb{R}\).
Comme \(\lim_{n \rightarrow + \infty} U_n = 0\) et que \(f(0) = 0\), 
la suite \((V_n)\) définie par \(V_n = f(U_n)\) est convergente et sa limite est \(0\).
\end{frame}
\begin{frame}{Suites numériques}
Soient \((U_n)\) et \((V_n)\) deux suites numériques telles que \(U_n \leq V_n\) à partir d'un certain rang.

\begin{itemize}
 
\item \(
      \text{Si }
      \lim_{n \rightarrow + \infty} U_n = + \infty
      \quad
      \text{ alors }
      \quad
      \lim_{n \rightarrow + \infty} V_n = + \infty
      \)

\item \(
      \text{Si }
      \lim_{n \rightarrow + \infty} V_n = - \infty
      \quad
      \text{ alors }
      \quad
      \lim_{n \rightarrow + \infty} U_n = - \infty
      \)

\end{itemize}
\end{frame}
\begin{frame}{Suites numériques}
Soient \((U_n)\) la suite définie par \(U_n = n^2 - 1 \) et \((V_n)\) la suite définie par \({V_n = n^2 + \sin (n)}\).
Comme \(-1 \leq \sin (n) \leq 1\) pour tout \(n \in \mathbb{N}\), alors 
      \(n^2 - 1 \leq n^2 + \sin (n) \leq n^2 + 1\) pour tout \(n \in \mathbb{N}\), ce qui veut dire que 
      \(U_n \leq V_n\) pour tout \(n \in \mathbb{N}\).
D'une part, on sait que \(\lim_{n \rightarrow + \infty} U_n = + \infty\); d'autre part, on sait que \(U_n \leq V_n\).
Ainsi, d'après le Théorème thm:comparaison suite, on a \(\lim_{n \rightarrow + \infty} V_n = + \infty\).
\end{frame}
\begin{frame}{Suites numériques}
Soient \((U_n)\), \((V_n)\) et \((W_n)\) trois suites numériques telles que \({U_n \leq V_n \leq W_n}\) à partir d'un certain rang.
Si \(\lim_{n \rightarrow + \infty} U_n = \lim_{n \rightarrow + \infty} W_n = \ell\) alors \(\lim_{n \rightarrow + \infty} V_n = \ell\).
\end{frame}
\begin{frame}{Suites numériques}
Considérons les trois suites \((U_n)\), \((V_n)\) et \((W_n)\) définies pour \(n \in \mathbb{N}^*\) respectivement par : 

\begin{itemize}
 
\item \( U_n = \dfrac{-1}{n} \)

\item \( V_n = \dfrac{\sin(n)}{n} \)

\item \( W_n = \dfrac{1}{n} \)

\end{itemize}
\end{frame}
\begin{frame}{Suites numériques (2)}
Comme \(-1 \leq \sin (n) \leq 1\) et \(n > 0\) pour tout \(n \in \mathbb{N}^*\), alors 
      \(U_n \leq V_n \leq W_n\).
Or, on sait que \( \lim_{n \rightarrow + \infty} U_n = \lim_{n \rightarrow + \infty} W_n = 0\); d'après le Théorème thm:gendarme suite, \( \lim_{n \rightarrow + \infty} V_n = 0 \).
\end{frame}
\begin{frame}{Suites numériques}
Suite du type \(a^n\)

 Soient \(a \in \mathbb{R}\) et \((U_n)\) la suite définie par \(U_n = a^n\) ; c'est la suite géométrique de raison \(a\) et de premier terme \(1\).
\end{frame}
\begin{frame}{Suites numériques (2)}
\begin{itemize}
 
\item Si \(a \leq -1\), la suite \((U_n)\) est divergente (elle n'a pas de limite).

\item Si \(a = 1\), la suite \((U_n)\) est constante, égale à \(1\).

\item Si \(-1 < a < 1\), la suite \((U_n)\) converge vers \(0\).

\item Si \(a > 1\), la suite \((U_n)\) tend vers \(+ \infty\).

\end{itemize}
\end{frame}
\begin{frame}{Suites numériques}
Suite du type \(n^a\)

 Soient \(a\) un nombre réel non nul et \((V_n)_{n \in \mathbb{N}^*}\) la suite numérique définie par \(V_n = n^a\)
(pour \(a = 0\), la suite est constante, égale à \(1\)).

\begin{itemize}
 
\item Si \(a < 0\), alors la suite \((V_n)\) converge vers \(0\).

\item Si \(a > 0\), alors la suite \((V_n)\) tend vers \(+\infty\).

\end{itemize}
\end{frame}
\begin{frame}{Suites numériques}
\textbf{Suite récurrente}

 Certaines suites peuvent être présentées sous forme de relation de récurrence et la donnée d'un premier terme.
Par exemple, considérons la suite \((U_n)\) définie par 
\[
\left\lbrace
\begin{array}{l}
U_0 = -2 \\
U_{n+1} = U_n + 3 .
\end{array}
\right.
\]

Comme \(U_{n+1} - U_n = 3\), \((U_n)\) est une suite arithmétique de raison \(3\).
Comme le premier terme est \(-2\), le terme général est \(U_n = -2 + 3n\).
Il n'est pas difficile de constater que la suite \((U_n)\) est croissante.
\end{frame}
\begin{frame}{Suites numériques (2)}
De même, les suites géométriques peuvent être présentées sous forme récurrente.
Considérons par exemple la suite \((V_n)\) définie par 
\[
\left\lbrace
\begin{array}{l}
V_0 = 2 \\
V_{n+1} = -3  V_n .
\end{array}
\right.
\]
\end{frame}
\begin{frame}{Suites numériques}
Comme \(\dfrac{V_{n+1}}{V_n} = -3\), \((V_n)\) est une suite géométrique de raison \(-3\).
Comme \(V_0 = 2\), l'expression de \(V_n\) en fonction de \(n\) est \(V_n = 2 \cdot (-3)^n\).
La suite \((V_n)\) est une suite alternée; en effet, le signe de \(V_n\) est différent du signe de \(V_{n+1}\).

Les suites arithmétiques et géométriques ne représentent que des cas très particulier de suites.
            Il existe des suites qui ne sont ni arithmétiques ni géométriques.
\end{frame}
\begin{frame}{Suites numériques (2)}
De même, les suites croissantes et décroissantes ne sont que des cas particuliers de suites.
            Il existe des suites qui ne sont ni croissante ni décroissante. 

Par contre, toute suite qui n'est pas convergente est divergente (c'est la définition).
\end{frame}
\begin{frame}{Suites numériques}
Suites définies par \(U_{n+1} = f(U_n)\) : représentation et limite

 Soit \(f\) une fonction définie sur un intervalle \(I\) telle que \(f(I) \subset I\), et \((U_n)\) définie par \(U_0 \in I\) et \(U_{n+1} = f(U_n)\).
\end{frame}
\begin{frame}{Suites numériques (2)}
\begin{itemize}

\item \emph{Représentation graphique.} On trace la courbe de \(f\) et la droite \(\Delta : y = x\).
      On place \(U_0\) sur l'axe des abscisses ; l'image \(U_1 = f(U_0)\) se lit sur l'axe des ordonnées et on la ramène
      sur l'axe des abscisses grâce à \(\Delta\) ; on recommence pour \(U_2\), \(U_3\), \ldots{}
      Ce dessin permet de conjecturer le sens de variation et la limite de la suite.

\end{itemize}
\end{frame}
\begin{frame}{Suites numériques}
Si la suite \((U_n)\) définie par \(U_{n+1} = f(U_n)\) converge vers un réel \(\ell\) appartenant à \(I\) et si \(f\) est continue en \(\ell\),
alors \(\ell\) vérifie \(f(\ell) = \ell\).
\end{frame}
\begin{frame}{Suites numériques}
Soit \((U_n)\) définie par \(U_0 = 0\) et \(U_{n+1} = \dfrac{1}{2} U_n + 3\).
On montre par récurrence que \(0 \leq U_n \leq U_{n+1} \leq 6\) : la suite est croissante et majorée, donc convergente.
Sa limite \(\ell\) vérifie \(\ell = \dfrac{1}{2}\ell + 3\), d'où \(\ell = 6\).
\end{frame}
\end{document}
